
\documentclass{article}
\usepackage{colortbl}
\usepackage{makecell}
\usepackage{multirow}
\usepackage{supertabular}

\begin{document}

\newcounter{utterance}

\twocolumn

{ \footnotesize  \setcounter{utterance}{1}
\setlength{\tabcolsep}{0pt}
\begin{supertabular}{c@{$\;$}|p{.15\linewidth}@{}p{.15\linewidth}p{.15\linewidth}p{.15\linewidth}p{.15\linewidth}p{.15\linewidth}}

    \# & $\;$A & \multicolumn{4}{c}{Game Master} & $\;\:$B\\
    \hline 

    \theutterance \stepcounter{utterance}  

    & & \multicolumn{4}{p{0.6\linewidth}}{\cellcolor[rgb]{0.9,0.9,0.9}{%
	\makecell[{{p{\linewidth}}}]{% 
	  \tt {\tiny [A$\langle$GM]}  
	 State giocando a un gioco di negoziazione in cui dovete accordarvi su come dividere una serie di oggetti tra voi e un altro giocatore.\\ \tt \\ \tt Regole:\\ \tt (a) A Lei e all'altro giocatore viene dato un insieme di oggetti. Ognuno di voi riceve anche una funzione di valore segreta, che rappresenta il valore di ciascun tipo di oggetto.\\ \tt (b) Si scambiano messaggi con l'altro giocatore per concordare chi si aggiudica gli oggetti. Potete inviare un massimo di 5 messaggi ciascuno, oppure terminare in anticipo facendo una proposta segreta in qualsiasi momento.\\ \tt (c) A ciascuno di voi viene chiesto di inviare una proposta segreta indicando gli oggetti che desiderate, formattata tra parentesi quadre come segue: "[Proposta: <numero> <nome oggetto>, <numero> <nome oggetto>, <...>]".\\ \tt (d) Se le vostre proposte sono complementari, cioè ci sono abbastanza oggetti per soddisfare entrambe le proposte, a ciascun giocatore viene assegnato un punteggio basato sulla somma dei valori degli oggetti ricevuti. In caso contrario, entrambi ricevono zero punti.\\ \tt (e) L'obiettivo è massimizzare il proprio punteggio.\\ \tt \\ \tt Cominciamo.\\ \tt \\ \tt L'insieme degli oggetti disponibili è:\\ \tt \\ \tt 1 arancia, 1 peperone, 3 quadri\\ \tt \\ \tt La funzione valore segreta è:\\ \tt \\ \tt arancia: 5, peperone: 2, quadro: 1\\ \tt \\ \tt IMPORTANTE: i vostri messaggi, a meno che non si tratti di una proposta segreta, vengono trasmessi direttamente all'altro giocatore, quindi non includete alcuna risposta alle regole o testo di annuncio del vostro messaggio. Per fare una proposta segreta, utilizzate il formato indicato. Non utilizzare delle parentesi quadre quando si comunica all'altro giocatore, altrimenti verrà interpretata come una proposta segreta.\\ \tt  
	  } 
	   } 
	   } 
	 & \\ 
 

    \theutterance \stepcounter{utterance}  

    & \multicolumn{4}{p{0.6\linewidth}}{\cellcolor[rgb]{0.8,1,0.9}{%
	\makecell[{{p{\linewidth}}}]{% 
	  \tt {\tiny [A$\rangle$GM]}  
	 Sono pronto a iniziare la negoziazione. Sembra che ci siano più quadri rispetto agli altri oggetti, quindi potrebbe essere una buona idea dividerli in modo equo. Che ne pensi di iniziare assegnando un paio di quadri a ciascuno di noi e poi discutere la divisione dell'arancia e del peperone? 
	  } 
	   } 
	   } 
	 & & \\ 
 

    \theutterance \stepcounter{utterance}  

    & & & \multicolumn{2}{p{0.3\linewidth}}{\cellcolor[rgb]{0.95,0.95,0.95}{%
	\makecell[{{p{\linewidth}}}]{% 
	  \tt {\tiny [GM$|$GM]}  
	 continue 
	  } 
	   } 
	   } 
	 & & \\ 
 

    \theutterance \stepcounter{utterance}  

    & & & \multicolumn{2}{p{0.3\linewidth}}{\cellcolor[rgb]{0.95,0.95,0.95}{%
	\makecell[{{p{\linewidth}}}]{% 
	  \tt {\tiny [GM$|$GM]}  
	 Sono pronto a iniziare la negoziazione. Sembra che ci siano più quadri rispetto agli altri oggetti, quindi potrebbe essere una buona idea dividerli in modo equo. Che ne pensi di iniziare assegnando un paio di quadri a ciascuno di noi e poi discutere la divisione dell'arancia e del peperone? 
	  } 
	   } 
	   } 
	 & & \\ 
 

    \theutterance \stepcounter{utterance}  

    & & \multicolumn{4}{p{0.6\linewidth}}{\cellcolor[rgb]{0.9,0.9,0.9}{%
	\makecell[{{p{\linewidth}}}]{% 
	  \tt {\tiny [GM$\rangle$B]}  
	 State giocando a un gioco di negoziazione in cui dovete accordarvi su come dividere una serie di oggetti tra voi e un altro giocatore.\\ \tt \\ \tt Regole:\\ \tt (a) A Lei e all'altro giocatore viene dato un insieme di oggetti. Ognuno di voi riceve anche una funzione di valore segreta, che rappresenta il valore di ciascun tipo di oggetto.\\ \tt (b) Si scambiano messaggi con l'altro giocatore per concordare chi si aggiudica gli oggetti. Potete inviare un massimo di 5 messaggi ciascuno, oppure terminare in anticipo facendo una proposta segreta in qualsiasi momento.\\ \tt (c) A ciascuno di voi viene chiesto di inviare una proposta segreta indicando gli oggetti che desiderate, formattata tra parentesi quadre come segue: "[Proposta: <numero> <nome oggetto>, <numero> <nome oggetto>, <...>]".\\ \tt (d) Se le vostre proposte sono complementari, cioè ci sono abbastanza oggetti per soddisfare entrambe le proposte, a ciascun giocatore viene assegnato un punteggio basato sulla somma dei valori degli oggetti ricevuti. In caso contrario, entrambi ricevono zero punti.\\ \tt (e) L'obiettivo è massimizzare il proprio punteggio.\\ \tt \\ \tt Cominciamo.\\ \tt \\ \tt L'insieme degli oggetti disponibili è:\\ \tt \\ \tt 1 arancia, 1 peperone, 3 quadri\\ \tt \\ \tt La funzione valore segreta è:\\ \tt \\ \tt arancia: 7, peperone: 3, quadro: 0\\ \tt \\ \tt IMPORTANTE: i vostri messaggi, a meno che non si tratti di una proposta segreta, vengono trasmessi direttamente all'altro giocatore, quindi non includete alcuna risposta alle regole o testo di annuncio del vostro messaggio. Per fare una proposta segreta, utilizzate il formato indicato. Non utilizzare delle parentesi quadre quando si comunica all'altro giocatore, altrimenti verrà interpretata come una proposta segreta.\\ \tt \\ \tt \\ \tt Sono pronto a iniziare la negoziazione. Sembra che ci siano più quadri rispetto agli altri oggetti, quindi potrebbe essere una buona idea dividerli in modo equo. Che ne pensi di iniziare assegnando un paio di quadri a ciascuno di noi e poi discutere la divisione dell'arancia e del peperone? 
	  } 
	   } 
	   } 
	 & \\ 
 

    \theutterance \stepcounter{utterance}  

    & & & \multicolumn{4}{p{0.6\linewidth}}{\cellcolor[rgb]{1,0.85,0.72}{%
	\makecell[{{p{\linewidth}}}]{% 
	  \tt {\tiny [GM$\langle$B]}  
	 Penso che dividere i quadri equamente sia una buona base di partenza, dato che ce ne sono tre. Potremmo effettivamente dare un paio di quadri a ciascuno di noi, lasciando un quadro da decidere. Per quanto riguarda l'arancia e il peperone, potremmo discuterne ulteriormente, visto che hanno valore per almeno uno di noi. Che ne pensi di scambiare informazioni sui nostri interessi per questi due oggetti, senza rivelare troppo sugli oggetti più preziosi per noi? 
	  } 
	   } 
	   } 
	  \\ 
 

    \theutterance \stepcounter{utterance}  

    & & & \multicolumn{2}{p{0.3\linewidth}}{\cellcolor[rgb]{0.95,0.95,0.95}{%
	\makecell[{{p{\linewidth}}}]{% 
	  \tt {\tiny [GM$|$GM]}  
	 continue 
	  } 
	   } 
	   } 
	 & & \\ 
 

    \theutterance \stepcounter{utterance}  

    & & & \multicolumn{2}{p{0.3\linewidth}}{\cellcolor[rgb]{0.95,0.95,0.95}{%
	\makecell[{{p{\linewidth}}}]{% 
	  \tt {\tiny [GM$|$GM]}  
	 Penso che dividere i quadri equamente sia una buona base di partenza, dato che ce ne sono tre. Potremmo effettivamente dare un paio di quadri a ciascuno di noi, lasciando un quadro da decidere. Per quanto riguarda l'arancia e il peperone, potremmo discuterne ulteriormente, visto che hanno valore per almeno uno di noi. Che ne pensi di scambiare informazioni sui nostri interessi per questi due oggetti, senza rivelare troppo sugli oggetti più preziosi per noi? 
	  } 
	   } 
	   } 
	 & & \\ 
 

    \theutterance \stepcounter{utterance}  

    & & \multicolumn{4}{p{0.6\linewidth}}{\cellcolor[rgb]{0.9,0.9,0.9}{%
	\makecell[{{p{\linewidth}}}]{% 
	  \tt {\tiny [A$\langle$GM]}  
	 Penso che dividere i quadri equamente sia una buona base di partenza, dato che ce ne sono tre. Potremmo effettivamente dare un paio di quadri a ciascuno di noi, lasciando un quadro da decidere. Per quanto riguarda l'arancia e il peperone, potremmo discuterne ulteriormente, visto che hanno valore per almeno uno di noi. Che ne pensi di scambiare informazioni sui nostri interessi per questi due oggetti, senza rivelare troppo sugli oggetti più preziosi per noi? 
	  } 
	   } 
	   } 
	 & \\ 
 

    \theutterance \stepcounter{utterance}  

    & \multicolumn{4}{p{0.6\linewidth}}{\cellcolor[rgb]{0.8,1,0.9}{%
	\makecell[{{p{\linewidth}}}]{% 
	  \tt {\tiny [A$\rangle$GM]}  
	 Sembra un approccio ragionevole. Dato che i quadri sono divisi in modo abbastanza equo, possiamo concentrarci sull'arancia e il peperone. Io sono interessato a uno di questi due oggetti, ma non sono disposto a rivelare quale per ora. Potresti dirmi se hai una preferenza per l'arancia o il peperone, o se sei disposto a considerare una proposta che includa entrambi? Il quadro rimanente potrebbe essere un buon oggetto di scambio per assicurarsi di ottenere l'altro oggetto desiderato. 
	  } 
	   } 
	   } 
	 & & \\ 
 

    \theutterance \stepcounter{utterance}  

    & & & \multicolumn{2}{p{0.3\linewidth}}{\cellcolor[rgb]{0.95,0.95,0.95}{%
	\makecell[{{p{\linewidth}}}]{% 
	  \tt {\tiny [GM$|$GM]}  
	 continue 
	  } 
	   } 
	   } 
	 & & \\ 
 

    \theutterance \stepcounter{utterance}  

    & & & \multicolumn{2}{p{0.3\linewidth}}{\cellcolor[rgb]{0.95,0.95,0.95}{%
	\makecell[{{p{\linewidth}}}]{% 
	  \tt {\tiny [GM$|$GM]}  
	 Sembra un approccio ragionevole. Dato che i quadri sono divisi in modo abbastanza equo, possiamo concentrarci sull'arancia e il peperone. Io sono interessato a uno di questi due oggetti, ma non sono disposto a rivelare quale per ora. Potresti dirmi se hai una preferenza per l'arancia o il peperone, o se sei disposto a considerare una proposta che includa entrambi? Il quadro rimanente potrebbe essere un buon oggetto di scambio per assicurarsi di ottenere l'altro oggetto desiderato. 
	  } 
	   } 
	   } 
	 & & \\ 
 

    \theutterance \stepcounter{utterance}  

    & & \multicolumn{4}{p{0.6\linewidth}}{\cellcolor[rgb]{0.9,0.9,0.9}{%
	\makecell[{{p{\linewidth}}}]{% 
	  \tt {\tiny [GM$\rangle$B]}  
	 Sembra un approccio ragionevole. Dato che i quadri sono divisi in modo abbastanza equo, possiamo concentrarci sull'arancia e il peperone. Io sono interessato a uno di questi due oggetti, ma non sono disposto a rivelare quale per ora. Potresti dirmi se hai una preferenza per l'arancia o il peperone, o se sei disposto a considerare una proposta che includa entrambi? Il quadro rimanente potrebbe essere un buon oggetto di scambio per assicurarsi di ottenere l'altro oggetto desiderato. 
	  } 
	   } 
	   } 
	 & \\ 
 

    \theutterance \stepcounter{utterance}  

    & & & \multicolumn{4}{p{0.6\linewidth}}{\cellcolor[rgb]{1,0.85,0.72}{%
	\makecell[{{p{\linewidth}}}]{% 
	  \tt {\tiny [GM$\langle$B]}  
	 Sono aperto a considerare varie possibilità per l'arancia e il peperone. Non ho una preferenza forte per uno_specifico tra i due, ma sarei disposto a fare una proposta che prenda in considerazione entrambi. Il quadro rimanente potrebbe essere usato come "moneta di scambio" per assicurarsi di ottenere l'oggetto più desiderato. Tuttavia, dovrei sapere un po' di più sulla tua disponibilità a cedere uno degli oggetti senza ottenere in cambio il quadro rimanente. Potresti essere disposto a cedere l'arancia o il peperone senza ricevere il terzo quadro in cambio? 
	  } 
	   } 
	   } 
	  \\ 
 

    \theutterance \stepcounter{utterance}  

    & & & \multicolumn{2}{p{0.3\linewidth}}{\cellcolor[rgb]{0.95,0.95,0.95}{%
	\makecell[{{p{\linewidth}}}]{% 
	  \tt {\tiny [GM$|$GM]}  
	 continue 
	  } 
	   } 
	   } 
	 & & \\ 
 

    \theutterance \stepcounter{utterance}  

    & & & \multicolumn{2}{p{0.3\linewidth}}{\cellcolor[rgb]{0.95,0.95,0.95}{%
	\makecell[{{p{\linewidth}}}]{% 
	  \tt {\tiny [GM$|$GM]}  
	 Sono aperto a considerare varie possibilità per l'arancia e il peperone. Non ho una preferenza forte per uno_specifico tra i due, ma sarei disposto a fare una proposta che prenda in considerazione entrambi. Il quadro rimanente potrebbe essere usato come "moneta di scambio" per assicurarsi di ottenere l'oggetto più desiderato. Tuttavia, dovrei sapere un po' di più sulla tua disponibilità a cedere uno degli oggetti senza ottenere in cambio il quadro rimanente. Potresti essere disposto a cedere l'arancia o il peperone senza ricevere il terzo quadro in cambio? 
	  } 
	   } 
	   } 
	 & & \\ 
 

    \theutterance \stepcounter{utterance}  

    & & \multicolumn{4}{p{0.6\linewidth}}{\cellcolor[rgb]{0.9,0.9,0.9}{%
	\makecell[{{p{\linewidth}}}]{% 
	  \tt {\tiny [A$\langle$GM]}  
	 Sono aperto a considerare varie possibilità per l'arancia e il peperone. Non ho una preferenza forte per uno_specifico tra i due, ma sarei disposto a fare una proposta che prenda in considerazione entrambi. Il quadro rimanente potrebbe essere usato come "moneta di scambio" per assicurarsi di ottenere l'oggetto più desiderato. Tuttavia, dovrei sapere un po' di più sulla tua disponibilità a cedere uno degli oggetti senza ottenere in cambio il quadro rimanente. Potresti essere disposto a cedere l'arancia o il peperone senza ricevere il terzo quadro in cambio? 
	  } 
	   } 
	   } 
	 & \\ 
 

    \theutterance \stepcounter{utterance}  

    & \multicolumn{4}{p{0.6\linewidth}}{\cellcolor[rgb]{0.8,1,0.9}{%
	\makecell[{{p{\linewidth}}}]{% 
	  \tt {\tiny [A$\rangle$GM]}  
	 Capisco il tuo punto di vista. Io sono disposto a considerare la cessione di uno degli oggetti senza ottenere il quadro rimanente in cambio, ma dipende molto da come vedi la situazione e da cosa sei disposto a offrire in cambio. Se potessi ottenere una buona sensazione di come valuti l'arancia e il peperone, potrei essere più aperto a fare una proposta che non include il quadro rimanente. In alternativa, se riuscissimo a trovare un accordo secondo cui ciascuno di noi ottenga l'oggetto che più gli interessa, potremmo dividere il quadro rimanente in modo che entrambi siamo soddisfatti del risultato finale. [Proposta: 2 quadri, arancia] 
	  } 
	   } 
	   } 
	 & & \\ 
 

    \theutterance \stepcounter{utterance}  

    & & & \multicolumn{2}{p{0.3\linewidth}}{\cellcolor[rgb]{0.95,0.95,0.95}{%
	\makecell[{{p{\linewidth}}}]{% 
	  \tt {\tiny [GM$|$GM]}  
	 proposal does not include number/name pairs 
	  } 
	   } 
	   } 
	 & & \\ 
 

\end{supertabular}
}

\end{document}
